\documentclass{beamer}
\usepackage[utf8]{inputenc}
\usetheme{Warsaw}


%Information to be included in the title page:
\title{DRASIL}
\subtitle{A Knowledge-Based Approach to Scientific Software Development}
\author{Aaron M, Dan S, Maryyam N, Nicholas R, Henry M}
\institute{McMaster University}
\date[LSS \today]{Literate Scientific Software Group, \today}


\begin{document}

\frame{\titlepage}

\begin{frame}
\frametitle{Background Context}
\begin{itemize}
 \item $\exists$ problems $\in$ D where
 \item $D = \{$ scientific computing, engineering computing $\}$
 \item Problems = [\\ 
  \begin{itemize}
   \item Inconsistent Software Requirement Specifications (SRS) across D
   \item Inconsistency between code and documentation
   \item Documentation is annoying to make and maintain
   \item Hard to reuse code for different applications
  \end{itemize}
 ]
\end{itemize}
\end{frame}

\begin{frame}
\frametitle{The Goal}
\includegraphics[scale=0.70]{../WG2_11/generate_all_the_things.jpg}
\end{frame}

\begin{frame}
\frametitle{Purpose of Drasil}
\begin{itemize}
  \item Solve the four issues
  \item Promote 
    \begin{itemize}
      \item Reusability 
        \begin{itemize}
          \item Examples have fully documented code
          \item Data base to build new examples
        \end{itemize}
      \item Maintainability
        \begin{itemize}
          \item Make changes in one place, gets updated everywhere
        \end{itemize}
    \end{itemize}
\end{itemize}
\end{frame}

\begin{frame}
\frametitle{What is Drasil?}
\begin{itemize}
  \item<1-> Knowledge Capture (Data.Drasil)
  \item<2-> Language and Rendering (Language.Drasil)
    \begin{itemize}
      \item Code Generation: transition from Drasil to working code
      \item Documentation Generation: transition from Drasil to human readable documentation
    \end{itemize}
  \item<3-> Case Studies (Example.Drasil)
    \begin{itemize}
      \item This part is where you would input equations, requirements, and output code and documentation
    \end{itemize}
\end{itemize}
\end{frame}

\begin{frame}
\frametitle{Introduction}
\begin{itemize}
 \item<1-> Scientific and engineering computing has the potential to lead other fields of software with its solid knowledge base
 \item<2-> Drasil is intended to simplify the generation of documentation and code for scientific software
 \item<3-> Facilitate desirable software qualities such as traceability, verifiability, and reproducibility
 \item<4-> Case studies from which structural patterns and implicit relationships can be extracted, data can be captured, and core systems can be tested and implemented
\end{itemize}
\end{frame}

\begin{frame}
\frametitle{}
\begin{itemize}
 \item<1-> Finding patterns within examples $\Rightarrow$ sentence combinators
 \item<2-> Finding patterns between examples $\Rightarrow$ extraction of common sections, contents, and concepts
 \item<3-> Knowledge extraction
 \item<4-> Reduce duplication
   \begin{itemize}
     \item Function efficiency
     \item Building chunks off of each other
   \end{itemize}
 \item<5-> Implement new functions/types created by supervisors
 \item<6-> Bug fixing
 \item<7-> Opening/closing issues
\end{itemize}
\end{frame}

\begin{frame}
\frametitle{Case Study Contributions}
\begin{itemize}
\item<1-> SWHS
  \begin{itemize}
    \item Largest Example
    \item ODEs
  \end{itemize}
\item<2-> NoPCM
  \begin{itemize}
    \item Builds of pre-existing SWHS example
  \end{itemize}
\item<3-> GlassBR
  \begin{itemize}
    \item General definition's omitted
  \end{itemize}
\item<4-> SSP
  \begin{itemize}
    \item Indexing
    \item Sophisticated math
    \item Diversity of symbols
  \end{itemize}
\item<5-> GamePhysics
  \begin{itemize}
    \item Most ambiguous example
    \item SRS for a game physics library
  \end{itemize}
\end{itemize}
\end{frame}

\begin{frame}
\frametitle{Collaboration via Github}
Git is a version control system, github is a git repository hosting service that is \alert{free}.
\begin{itemize}
 \item<1-> Git allows us to collaborate effectively, even when team members are not in the same location
 \item<2-> Git combined with haskell, allows us to make large changes while easily maintaining a working version of Drasil
 \item<3-> Git \alert{(when used properly)} prevents catastrophic lose of work
\end{itemize}
\end{frame}

\begin{frame}
\frametitle{End}
For more information about Drasil and LLS visit our github page: \\
https://github.com/JacquesCarette/literate-scientific-software \\
\alert{You can even build a working version yourself!}
\end{frame}

\end{document}